%% ---------------------------------------------------------------------------
%% bilingual-template.tex
%%
%% A "parallel text" edition: the original on the left, the translation on the
%% right, with paragraph N always facing paragraph N - across page breaks too.
%%
%% BUILD:  ./build.sh      (or: pdflatex bilingual-template.tex, twice)
%%
%% HOW TO WRITE YOUR OWN TEXT
%%   Everything you type lives inside \begin{paracol}{2} ... \end{paracol}.
%%   One paragraph = one \pair command:
%%
%%       \pair{Original paragraph}{Translated paragraph}
%%
%%   First argument  -> left column,  typeset in ORIGLANG
%%   Second argument -> right column, typeset in TRANSLANG
%%
%%   To switch languages, edit \origlang and \translang near the top of this
%%   file. Braces may span several lines, so feel free to wrap long paragraphs.
%%
%% HOW THE ALIGNMENT WORKS (this is the whole trick)
%%   \pair typesets each side into a box, works out which one is TALLER, and
%%   pads the shorter one with invisible space so that both boxes come out at
%%   exactly the same height. Both columns therefore advance by an identical
%%   amount after every paragraph, which is what keeps paragraph N facing
%%   paragraph N - and what makes the page break fall in the same place in
%%   both languages.
%%
%%   Note: paracol's own \switchcolumn does NOT do this. It alternates columns
%%   but lets the two texts drift apart by a line whenever a translation runs
%%   longer, so the lockstep is done here instead. \switchcolumn is used purely
%%   to move between the columns.
%%
%%   Consequence: a paragraph is never split across a page. Keep them shorter
%%   than the text height (see \textheight below) or you will get a warning.
%% ---------------------------------------------------------------------------
\documentclass[12pt,a4paper,twoside]{book}

%% --- the two languages of this edition -------------------------------------
%% Available ones you can drop in: english, french, spanish, german, italian,
%% portuguese, polish, czech, latin... (any babel-supported language)
\newcommand{\origlang}{french}
\newcommand{\translang}{spanish}

%% --- encoding & hyphenation -------------------------------------------------
\usepackage[T1]{fontenc}
\usepackage[utf8]{inputenc}
\usepackage[\origlang,\translang]{babel}

%% --- page geometry ---------------------------------------------------------
%% "inner"/"outer" (rather than left/right) is what you want for a book: the
%% wide margin falls on the binding side whichever page we are on.
\usepackage[inner=30mm,outer=22mm,top=25mm,bottom=25mm]{geometry}

%% --- typography ------------------------------------------------------------
\usepackage{lmodern}
\usepackage{microtype}

%% --- two columns -----------------------------------------------------------
\usepackage{paracol}

\columnratio{0.5,0.5}
\setlength{\columnsep}{9mm}
\setlength{\columnseprule}{0.3pt}   % hairline between the two texts
\setlength{\parindent}{1em}         % indent paragraphs, no extra space
\setlength{\parskip}{0pt}

%% --- running heads ---------------------------------------------------------
%% NB: in fancyhdr, L/R = physical side of the page, O/E = odd/even page.
%% The outer edge of a recto page is on the RIGHT and of a verso page on the
%% LEFT, so these have to be [RO] and [LE] - not [LO]/[RE], which would put the
%% heads on the binding edge.
\usepackage{fancyhdr}
\pagestyle{fancy}
\setlength{\headheight}{14.5pt}   % fancyhdr needs room for \small\itshape
\fancyhf{}
\fancyhead[RO]{\small\itshape Le réverbère du numéro douze}
\fancyhead[LE]{\small\itshape El farol del número doce}
\fancyfoot[LE,RO]{\small\thepage}
\renewcommand{\headrulewidth}{0.4pt}
\renewcommand{\footrulewidth}{0pt}

%% --- a heading that behaves inside a column --------------------------------
%% (\section would try to start a new float/page, so we do it by hand)
\newcommand{\heading}[1]{%
  \par{\large\bfseries #1\par}\nobreak\vspace{0.4\baselineskip}}

%% =========================== THE ALIGNMENT ENGINE ==========================
\newsavebox{\pairA}
\newsavebox{\pairB}
\newlength{\pairht}
\newlength{\pairgap}

%% Space left between two paragraphs.  Because \nointerlineskip is used below,
%% this (and not \baselineskip) is what separates one paragraph from the next.
%% ~0.35\baselineskip reproduces the usual book spacing.  It can only be
%% computed once \baselineskip exists, hence \AtBeginDocument.
\AtBeginDocument{\setlength{\pairgap}{0.35\baselineskip}}

%% Put one measured paragraph into the current column, stretched to \pairht so
%% that it occupies exactly the same vertical space in both columns.
\newcommand{\pairplace}[1]{%
  \par\nointerlineskip
  \noindent\vbox to\pairht{\unvbox#1\vfil}\par
  \nointerlineskip\vskip\pairgap\nointerlineskip
}

%% --- THE ONE COMMAND YOU NEED ---------------------------------------------
%% Original first, translation second. Each side gets its own language.
\newcommand{\pair}[2]{%
  \par
  \setbox\pairA=\vtop{\hsize=\linewidth\selectlanguage{\origlang}#1\relax}%
  \setbox\pairB=\vtop{\hsize=\linewidth\selectlanguage{\translang}#2\relax}%
  %% the slot is as tall as the taller of the two paragraphs
  \pairht=\ht\pairA \advance\pairht by \dp\pairA
  \dimen0=\ht\pairB \advance\dimen0 by \dp\pairB
  \ifdim\dimen0>\pairht \pairht=\dimen0 \fi
  \ifdim\pairht>\textheight
    \PackageWarning{bilingual}{%
      A paragraph is \the\pairht\space tall but \textheight\space is the page
      text height. It cannot be split, so it will overflow the page. Split it
      into two \pair commands.}%
  \fi
  \pairplace{\pairA}%
  \switchcolumn
  \pairplace{\pairB}%
  \switchcolumn}

\begin{document}

\begin{paracol}{2}

\pair
  {\heading{Chapitre premier}}
  {\heading{Capítulo primero}}

%% ============================ SAMPLE TEXT ==================================
%% 10 placeholder paragraphs of original prose, written to be roughly the same
%% length in both languages.  Delete them and put your own text in their place.
%% They run over onto further pages on purpose: check that each page starts the
%% same paragraph at the same height in both columns.

\pair
  {\noindent Dans une ville que les géographes avaient oubliée, un vieux
   réverbère s'était mis à clignoter trois soirs de suite. Personne n'osa
   l'éteindre, car on disait qu'il guidait encore un marinier perdu. À cette
   époque, les capitaines ne réclamaient plus de fanal au port : ils y
   arrivaient en comptant les tours de la ville. Mais depuis que le réverbère
   clignotait, les enfants du quartier s'étaient mis à compter les éclats
   eux-mêmes, et ils prétendaient que le nombre variait selon l'humeur de la mer.}
  {\noindent En una ciudad que los geógrafos habían olvidado, un viejo farol
   había empezado a parpadear tres noches seguidas. Nadie se atrevió a
   apagarlo, porque se decía que todavía guiaba a un marinero perdido. En
   aquel tiempo los capitanes ya no pedían un faro en el puerto: llegaban
   guiándose por las torres de la ciudad. Pero desde que el farol parpadeaba,
   los niños del barrio se habían puesto a contar los destellos ellos mismos,
   y decían que el número variaba según el humor del mar.}

\pair
  {Le garçon qui habitait au numéro douze passa devant la boutique fermée et
   frappa trois coups, comme on frappe à une porte que l'on connaît. Il n'y
   avait personne derrière, mais la boutique du réverbère ne fermait jamais
   tout à fait : la vendeuse laissait la clé de la cave sur le rebord de la
   fenêtre, par habitude, comme on laisse une lampe allumée pour celui qui
   rentre tard.}
  {El niño que vivía en el número doce pasó por delante de la tienda cerrada
   y golpeó tres veces, como se golpea a una puerta que se conoce. No había
   nadie detrás, pero la tienda del farol no cerraba nunca del todo: la
   dependienta dejaba la llave del sótano en el alféizar de la ventana, por
   costumbre, como se deja una lámpara encendida para quien llega tarde.}

\pair
  {Il attendit longtemps, la pluie fine sur les épaules, et finit par
   s'asseoir sur le pavé, les genoux serrés contre la poitrine. Le réverbère
   clignotait par-dessus les toits, régulier comme une pendule, et chaque
   éclat posait une petite lumière bleue sur les tuiles, une lumière qui
   s'éteignait avant d'avoir fini de briller. Il se demanda combien de fois
   il faudrait qu'il clignote pour qu'un bateau se décide à rentrer.}
  {Esperó mucho tiempo, con la llovizna fina sobre los hombros, y acabó
   sentándose en el adoquín, con las rodillas apretadas contra el pecho. El
   farol parpadeaba por encima de los tejados, regular como un reloj, y cada
   destello ponía una pequeña luz azul sobre las tejas, una luz que se
   apagaba antes de haber terminado de brillar. Se preguntó cuántas veces
   tendría que parpadear para que un barco se decidiera a volver.}

\pair
  {La vendeuse parut enfin, tenant un chiffon dans une main et la clé de la
   cave dans l'autre. Elle le regarda longtemps sans rien dire, comme on mesure
   une chose que l'on n'a pas l'habitude de voir. Puis elle demanda s'il
   savait monter six cent marches. Il répondit qu'il était capable d'en monter
   mille, et qu'il l'avait déjà fait le jour où il avait perdu son
   laissez-passer sous un banc. Elle ne rit pas, mais elle cessa de regarder
   la pluie.}
  {Por fin apareció la dependienta, con un trapo en una mano y la llave del
   sótano en la otra. Lo miró mucho rato sin decir nada, como se mide una cosa
   que no se está acostumbrado a ver. Luego preguntó si sabía subir seiscientos
   peldaños. Él respondió que era capaz de subir mil, y que ya lo había hecho
   el día en que perdió su pase de navegación debajo de un banco. Ella no
   rio, pero dejó de mirar la lluvia.}

\pair
  {« Il faut que quelqu'un monte là-haut, dit-elle enfin. Le réverbère doit
   briller cette nuit, sinon le marinier ne rentrera pas. » Elle expliqua
   qu'un réverbère qui cesse de briller n'éteint pas seulement une lumière :
   il retire aux bateaux la ligne qu'ils suivent depuis toujours. Le garçon
   demanda alors si le marinier existait encore. Elle dit qu'elle n'en était
   pas sûre, et que c'était précisément pour cela qu'il ne fallait pas laisser
   la lumière mourir.}
  {«Es preciso que alguien suba ahí arriba —dijo por fin—. El farol tiene que
   brillar esta noche, o el marinero no volverá.» Explicó que un farol que
   deja de brillar no apaga solamente una luz: le quita a los barcos la línea
   que siguen desde siempre. El niño preguntó entonces si el marinero existía
   todavía. Ella dijo que no estaba segura, y que precisamente por eso no
   había que dejar morir la luz.}

\pair
  {Le garçon ne demanda pas qui était le marinier. Il prit simplement la
   clé, et monta les six cent marches sans les compter, parce que les nombres
   lui paraissaient alors une invention des adultes. Il monta, et la ville
   changea de forme à mesure qu'il montait : d'abord les ruelles, puis les
   toits, puis les cheminées, puis la mer, immense et plate, qui n'avait plus
   l'air d'une chose que l'on pût traverser.}
  {El niño no preguntó quién era el marinero. Tomó simplemente la llave y
   subió los seiscientos peldaños sin contarlos, porque los números le
   parecían entonces una invención de los adultos. Subió, y la ciudad cambió
   de forma a medida que ascendía: primero las callejuelas, luego los tejados,
   luego las chimeneas, y por fin el mar, inmenso y llano, que ya no tenía el
   aspecto de una cosa que se pudiera atravesar.}

\pair
  {En haut, la ville s'étendait, mince et grise, et l'on distinguait à peine
   la mer. Le réverbère attendait, noir et silencieux, comme une bouche qui
   aurait oublié ses paroles. Sa lanterne était fendue, son cuivre verdi, et
   les pigeons du clocher s'étaient rassemblés sur la galerie comme s'ils
   savaient ce qui allait se passer. Il y avait de la poussière sur la marche
   supérieure, et personne n'avait nettoyé la place depuis des années.}
  {Arriba, la ciudad se extendía, delgada y gris, y apenas se distinguía el
   mar. El farol esperaba, negro y silencioso, como una boca que hubiera
   olvidado sus palabras. Su linterna estaba rajada, el cobre se había puesto
   verde, y las palomas de la campana se habían juntado en la galería como si
   supieran lo que iba a pasar. Había polvo en el peldaño de arriba, y nadie
   había limpiado aquel sitio en años.}

\pair
  {Il monta encore quatre marches et trouva, dans la loge de fer, une boîte de
   lampes de rechange et un carnet couvert de toile cirée. Le carnet
   appartenait à l'homme qui avait construit le réverbère, cinquante ans plus
   tôt. Il l'ouvrit à la dernière page et trouva une phrase écrite d'une main
   d'enfant : « N'oublie pas de regarder en haut. » Au-dessous, d'une autre
   main, une réponse : « Je n'oublie jamais. »}
  {Subió otros cuatro peldaños y encontró, en la caseta de hierro, una caja de
   lámparas de repuesto y un cuaderno con tapas de hule. El cuaderno había
   sido del hombre que construyó el farol, cincuenta años antes. Lo abrió por
   la última página y encontró una frase escrita con letra de niño: «No
   olvides de mirar hacia arriba.» Debajo, con otra letra, una respuesta:
   «No lo olvido nunca.»}

\pair
  {Il alluma la lampe avec la dernière allumette et s'assit sur la galerie. La
   lumière s'étala sur les toits, sur les cheminées, sur la pointe des mâts,
   et la mer cessa d'être une chose sombre. Le réverbère, qui n'avait pas
   servi depuis cinquante ans, se remit à clignoter comme s'il n'avait jamais
   cessé. En bas, la ville reprit son travail : on entendit les volets, les
   pas, une porte, puis une autre porte.}
  {Encendió la lámpara con la última cerilla y se sentó en la galería. La luz
   se derramó sobre los tejados, sobre las chimeneas, sobre las puntas de los
   mástiles, y el mar dejó de ser una cosa oscura. El farol, que no servía
   desde hacía cincuenta años, empezó otra vez a parpadear como si nunca
   hubiera parado. Abajo, la ciudad retomó su trabajo: se oyeron las
   contraventanas, los pasos, una puerta y luego otra puerta.}

\pair
  {Le marinier rentra à minuit. Il leva les yeux vers la fenêtre de la tour et
   dit merci à voix basse, comme on remercie une ville qui ne demande rien.
   Le garçon, resté sur la galerie, entendit le bateau se ranger contre le
   quai, entendit une corde, entendit une voix qui fredonnait. Il dormit là
   jusqu'au matin, et se réveilla avec la lumière du réverbère qui
   s'éteignait.}
  {El marinero volvió a medianoche. Levantó los ojos hacia la ventana de la
   torre y dio las gracias en voz baja, como se agradece a una ciudad que no
   pide nada. El niño, que se había quedado en la galería, oyó el barco
   atracar en el muelle, oyó una cuerda, oyó una voz que tarareaba. Durmió
   allí hasta la mañana y se despertó con la luz del farol apagándose.}

%% End of the text. \end{paracol} takes care of the last column.
\end{paracol}

\end{document}
